\documentclass[11pt]{report}
\usepackage[utf8]{inputenc}
\usepackage[T1]{fontenc}
\usepackage{lmodern}
\usepackage{amsmath,amssymb}
\usepackage{graphicx}
\usepackage[margin=1in]{geometry}
\usepackage{hyperref}
\usepackage{booktabs}
\usepackage[numbers]{natbib}
\usepackage{xcolor}
\usepackage{tcolorbox}
\usepackage{siunitx}
\usepackage{float}

% Colors
\definecolor{sectionblue}{HTML}{1e40af}
\definecolor{highlightbg}{HTML}{eff6ff}
\definecolor{notebg}{HTML}{fef3c7}

% Highlight box
\newtcolorbox{highlight}{
  colback=highlightbg,
  colframe=sectionblue!50,
  boxrule=0.5pt,
  arc=3pt,
  left=6pt, right=6pt, top=6pt, bottom=6pt
}

% Note box
\newtcolorbox{note}{
  colback=notebg,
  colframe=orange!50,
  boxrule=0.5pt,
  arc=3pt,
  left=6pt, right=6pt, top=6pt, bottom=6pt
}

\hypersetup{
  colorlinks=true,
  linkcolor=sectionblue,
  citecolor=sectionblue,
  urlcolor=sectionblue
}

\title{\textbf{Scientific Report Title}\\[0.5em]
\large Subtitle or Project Name}
\author{Author Name\\
\small Institution or Laboratory\\
\small \href{mailto:author@institution.edu}{author@institution.edu}}
\date{\today}

\begin{document}

\maketitle

\begin{abstract}
Provide a concise summary of the research objectives, methodology, key findings, and conclusions. The abstract should be self-contained and typically 150--300 words.
\end{abstract}

\tableofcontents

\chapter{Introduction}

Describe the background, motivation, and objectives of the study. Include relevant literature context and clearly state the research questions or hypotheses.

\begin{highlight}
\textbf{Research Question:} State your primary research question or hypothesis here.
\end{highlight}

\chapter{Materials and Methods}

\section{Experimental Design}

Describe the overall experimental approach, including study design, variables, and controls.

\section{Data Collection}

Detail the instruments, protocols, and procedures used for data collection.

\section{Statistical Analysis}

All analyses were performed using standard statistical methods. Results were considered significant at $p < 0.05$.

\begin{note}
\textbf{Note:} Include details about software versions, packages, and specific parameters used in the analysis.
\end{note}

\chapter{Results}

\section{Descriptive Statistics}

Present summary statistics and initial findings.

\begin{table}[H]
\centering
\caption{Summary statistics of key variables.}
\label{tab:summary}
\begin{tabular}{lrrr}
\toprule
\textbf{Variable} & \textbf{Mean} & \textbf{SD} & \textbf{N} \\
\midrule
Variable A & \num{42.5} & \num{3.2} & 100 \\
Variable B & \num{18.7} & \num{2.1} & 100 \\
Variable C & \num{7.3}  & \num{1.8} & 100 \\
\bottomrule
\end{tabular}
\end{table}

\section{Main Findings}

Present the main results with appropriate figures and tables. Reference Table~\ref{tab:summary} and any figures as needed.

\chapter{Discussion}

Interpret the results in the context of existing literature. Discuss implications, limitations, and future directions.

\chapter{Conclusion}

Summarize the key findings and their significance. State the main contributions of this work.

\bibliographystyle{plainnat}
\bibliography{references}

\end{document}
